%-----------------------------------------------------------------------------%
\chapter*{\Judul}
%-----------------------------------------------------------------------------%
\singlespacing
\begin{center}
    
    \vspace{-4em}
    
    \penulis
    
	\bigskip
    
    \textbf{ABSTRAK}
    
\end{center}

% \chapter*{Abstrak}

\vspace*{0.2cm}
{
	\setlength{\parindent}{0pt}

	\bigskip
	\bigskip

PT Pratama Nusantara Sakti merupakan perusahaan agrikultur perkebunan tebu yang
mengelola tenaga kerja kontrak dalam jumlah besar, namun masih mengandalkan
aplikasi berbasis Microsoft Access yang rentan terhadap kesalahan, duplikasi
data, dan celah kecurangan administratif. Kerja magang ini bertujuan untuk
merancang dan membangun \textit{HR Safety Campus}, yaitu sistem administrasi
tenaga kerja kontrak berbasis web yang dikembangkan menggunakan kerangka kerja
Laravel dengan arsitektur \textit{Model-View-Controller} (MVC) dan metode
pengembangan \textit{Waterfall}. Sistem ini mengimplementasikan arsitektur
\textit{Role-Based Access Control} (RBAC) untuk membatasi hak akses pengguna,
antarmuka responsif berbasis \textit{Tailwind CSS}, serta mekanisme penanganan
duplikasi data pada proses migrasi sekitar 4.000 rekaman tenaga kerja dari sistem
lama ke basis data MySQL. Hasil kerja magang berupa enam fitur utama, yaitu
registrasi tenaga kerja, klaim transportasi, surat izin keluar, manajemen tarif,
persetujuan karyawan, dan manajemen pengguna, yang berhasil diintegrasikan dan
diuji. Dengan adanya sistem ini, proses administrasi tenaga kerja menjadi lebih
cepat, akurat, dan terpusat, sekaligus menutup celah manipulasi data yang
sebelumnya merugikan perusahaan.
	\bigskip
 
% Kata kunci urut abjad
% 3 – 5 kata kunci
	\textbf{Kata kunci:} HRIS, Laravel, migrasi data, RBAC, tenaga kerja kontrak
} 

\onehalfspacing
